% =====================================================================
% 5.2 학술적 기여 — 5단계 5-2 (clean subsection wrapper)
% 본문 텍스트 무수정. 분할 본문 파일을 새 subsection 구조로 결합.
%   5.2.1 이론적 기여  ← 5.1.2 + 5.2(전체) + 5.8.1 본문
%   5.2.2 방법론적 기여 ← 5.3(전체) + 5.8.2 본문
% =====================================================================

\section{\vrev{학술적 기여}}
\label{sec:academic-contrib}

\begin{p5add}
본 학술적 기여는 본 연구가 BCM\chg{과 재해 복구}(DR) 학문 분야의
지식 체계에 새롭게 기여한 바를 \chg{\m{(i)} 이론적, \m{(ii)} 방법론적, 그리고 \m{(iii)} 실증적} 세 차원에서
종합한다. 이론적 기여(5.2.1항)에서는 $\DRTO$ 모델의 수학적 구조와
이중 채널 감쇠 메커니즘의 신규성을, 방법론적 기여(5.2.2항)에서는 비중첩
대역 시드 체계와 분할 회귀 방법론을, 실증적 기여(5.2.3항)에서는 \jrev{10개
가설} 전원 채택과 다중 검증 프레임워크의 결과\jrev{, 그리고 \textbf{분포 환경
적합성 종합 명제}(\secref{sec:ch4_hypothesis_summary})}를 다룬다.
\end{p5add}

% --- 5.2.1 이론적 기여 ---
\subsection{\vrev{이론적 기여}}
\label{subsec:academic-theoretical}

\begingroup
  % 원본 파일의 \section, \subsection 헤더만 시각적으로 억제 (label은 유지)
  \renewcommand{\section}[1]{}
  \renewcommand{\subsection}[1]{}
  \renewcommand{\subsubsection}[1]{\paragraph{#1}}

  % --- 주축: 5.1.2 이론적 기여 (간결한 intro + 1)/2)/3) 형식) ---
  %% -------------------------------------------------------------------
\subsection{\vrev{이론적 기여}}
\label{subsec:ch5_theoretical}
%% -------------------------------------------------------------------

본 연구의 이론적 기여는 네 가지 차원에서 정리된다.

\chg{첫째, 개별 시그모이드 바운딩에 의한 물리적 경계 보장이다.}
앞서 비교한 선형 RTO 모델의 물리적 경계 위반 문제를 개별 시그모이드
바운딩 함수로 해결하였다. $E_{O,\text{bnd}} \in (-R_0, 0]$과
$E_{U,\text{bnd}} \in [0, R_{\max} - R_0)$의 범위가 해석적으로 보장되어,
$\DRTO \in (0, R_{\max})$가 입력값에 무관하게 성립한다.
이는 기존 BCM 문헌에서 제시되지 않았던 수학적 보장으로서,
\chg{복구목표시간}의 음수 산출이나 MTPD 초과라는 구조적 문제를
근원적으로 해결한다.

\chg{둘째, 이중채널 감쇠 효과의 발견과 정량화이다.}
P85.8 교차점을 기준으로 $k_O$ 계수가 $0.52$배로 감쇠되는
이중채널 현상의 발견은 BCM 이론에 새로운 통찰을 제공한다.
이 발견은 극단 불확실성과 관측성이 비선형적으로 상호 작용하는
메커니즘을 규명하며, 시그모이드 비선형성에 의한 이론적 설명을
제시한다.

\chg{셋째, \citet{lee2026drto}과 본 연구의 일관적 이론 체계 구축이다.}
선행 연구\citep{lee2026drto}의 선형 근사$\to$시그모이드 바운딩$\to$$\kappa$ 최적화에서
본 연구의 이중채널 확장(C2, C3)까지의 단계적 확장이
\chg{수식·데이터·결론}의 교차 정합성을 유지하며 수행되었다.
선형 모델이 시그모이드 바운딩의 1차 근사임을 보이는 관계
($|z| \ll 1$에서 $S(z) \approx z$)는 이론적 확장의
내적 일관성을 수학적으로 보장한다.

  넷째, 불확실성의 확률적 정량화이다. 복구 과정의 불확실성을 일상적 변동을 나타내는 가우시안과 극단 사건을 나타내는 파레토의 두 분포로 정량화함으로써, 관측성 단일 차원에 머물던 기존 모델의 한계를 보완하고 극단 영역의 꼬리 위험을 복구목표시간 산정에 반영하였다.



  % --- 보완: 5.2 이론적 시사점 (기술적 세부 설명 보완) ---
  \section{\vrev{이론적 시사점}}
\label{sec:theoretical-implications}
%% ===================================================================

본 연구의 이론적 기여는 BCM 분야에 \jrev{\chg{i) 관측성 개념의 통합, ii) 불확실성의 확률적 정량화(가우시안과 파레토), iii) 개별 시그모이드 바운딩 이론의 정립, iv) 이중 채널 감쇠 효과의 발견이라는} 네 가지 축}으로 정리된다.
\p{[}표~\vref{tab:ch6-theoretical-contributions}\p{]}는 이론적 기여의 핵심 내용을 종합한다.

\begin{table}[htbp]
\centering
\caption{\jrev{이론적 기여 종합}}
\label{tab:ch6-theoretical-contributions}
\small
\begin{tabularx}{\textwidth}{c l X c}
\toprule
\textbf{\#} & \textbf{기여 영역} & \textbf{핵심 내용} & \textbf{근거} \\
\midrule
1 & 관측성 통합 &
  제어 이론의 관측성 개념을 BCM에 최초 도입;
  $z_O = \kappa \cdot k_O \cdot O_{\mathrm{raw}}$으로 정량화하여
  정적 RTO의 구조적 한계를 극복 &
  H1, \jrev{H3}-a \\
\addlinespace
2 & 시그모이드 바운딩 이론 &
  수정 시그모이드 $S(z) = 2/(1+e^{-z})-1$을 기저 함수로 채택;
  $\DRTO = R_0 + E_{O,bnd} + E_{U,bnd}$로 물리적 경계
  $(0,\; R_{\max})$를 해석적으로 보장 &
  \jrev{H2 ($\kappa = 1.2$ 전제)} \\
\addlinespace
3 & 관측성--불확실성 이원 구조 &
  관측성 효과($E_{O,bnd} \leq 0$)와 불확실성 효과($E_{U,bnd} \geq 0$)의
  구조적 반대 방향 확인; Cohen's $d$: $-1.180$ vs $+1.871$ &
  \jrev{H1, H8} \\
\addlinespace
4 & P85.8 CDF 교차점 규명 &
  가우시안(C2)과 파레토(C3) 분포의 $\eta$ 분위수 함수가
  P85.8에서 교차; 꼬리 효과의 조건부 특성 규명 &
  \jrev{H3}-c \\
\addlinespace
5 & 이중채널 감쇠 발견 &
  P85.8 기준 분할 회귀에서 Tail 영역의 $k_O$ 계수($-0.415$)가
  Core($-0.798$) 대비 $0.52$배 감쇠; 극단 환경에서 관측성
  효과의 상대적 축소 &
  \jrev{H8} \\
\addlinespace
6 & 꼬리 리스크 비선형 증폭 &
  CVaR$_{99}$: C3가 C2 대비 $+24.0\%$;
  분위수 상승에 따른 비선형 격차 확대
  (CVaR$_{90}$: $+20.4\%$ $\to$ CVaR$_{99}$: $+24.0\%$) &
  \jrev{H8} \\
\bottomrule
\end{tabularx}
\end{table}


%% -------------------------------------------------------------------
\subsection{\vrev{\revdel{BCM 이론에 대한 관측성 통합}}}
\label{subsec:observability-integration}
%% -------------------------------------------------------------------

\fdel{전통적 BCM 이론에서 RTO는 BIA(업무영향분석) 시점에 고정되는
정적 매개변수로 취급되어, 시스템 상태의 실시간 변화를 반영하지 못하였다.
본 연구는 제어 이론~\citep{kalman1960control}의 관측성(observability) 개념을
BCM 분야에 최초로 도입하여, 시스템의 내부 상태를 외부 출력으로부터
추정할 수 있는 정도를 $O_{\mathrm{raw}} \in [0,1]$로 정규화하고,
이를 관측성 가중치 $k_O$와 곡률 매개변수 $\kappa = 1.2$를 통해
$\DRTO$의 관측성 조정량으로 변환하였다.}

\fdel{관측성 조정량 $E_{O,bnd} = -R_0 \cdot S(\kappa \cdot k_O \cdot O_{\mathrm{raw}})$는
항상 비양수($\leq 0$)이므로, 관측성이 높을수록 $\DRTO$가 $R_0$ 이하로 단축된다.
이 구조는 \jrev{H3}-a에서 $100\%$ 준수율로 검증되었으며, 이는 확률적 결과가 아닌
모델의 \textbf{구조적 필연}이다. Cohen's $d = -1.180$의 매우 큰 효과 크기는
관측성 통합이 BCM의 복구 예측 정확도를 질적으로 향상시킨다는 이론적 근거를 제공한다.}

\fdel{이러한 관측성 통합은 기존 BCM 이론의 세 가지 공백을 해소한다.
첫째, RTO를 ``고정 상수''에서 ``상태 의존 함수''로 전환함으로써
동적 환경에서의 BCM 의사결정 기반을 확보한다.
둘째, APM(Application Performance Monitoring), SIEM(Security Information and
Event Management) 등 기존 모니터링 인프라의 BCM 활용 경로를 이론적으로 정당화한다.
이는 SRE(Site Reliability Engineering) 분야에서 관측성을 에러 예산(error budget)과
연계하여 서비스 신뢰성을 정량적으로 관리하는 접근(Beyer et al.\ 2016)과
맥락을 공유하며, DevOps 문화에서 모니터링 인프라 투자가 복구 성능을 향상시킨다는
실증적 발견(Forsgren et al.\ 2018; Kim et al.\ 2016)과도 부합한다.
셋째, 관측성 투자($k_O$ 증가)와 불확실성 허용($k_U$ 증가) 간의
trade-off를 $k_O + k_U = 1$ 제약으로 형식화하여,
BCM 자원 배분의 분석적 프레임워크를 제시한다.}


%% -------------------------------------------------------------------
\subsection{\vrev{\revdel{개별 시그모이드 바운딩 이론}}}
\label{subsec:sigmoid-bounding-theory}
%% -------------------------------------------------------------------

\fdel{$\DRTO$ 모델의 수학적 핵심은 수정 시그모이드 함수
$S(z) = 2/(1+e^{-z})-1 \in (-1,1)$에 의한 개별 바운딩(individual bounding)이다.
관측성 효과와 불확실성 효과를 각각 독립적인 시그모이드로 바운딩하여
$\DRTO = R_0 + E_{O,bnd} + E_{U,bnd}$로 정식화함으로써, 다음의 세 가지
이론적 요구사항을 동시에 충족한다.}

% [5단계 Plan H 재편: 구 5.2.2 "개별 시그모이드 바운딩 이론" 3개 항목
%  (물리적 경계 보장 / 무한 미분 가능성 / 교호작용 포착)은 5.2.1 이론적
%  기여로 통합·흡수됨. 기존 \fdel{} 래핑된 본문은 빈 enumerate 번호(1/2/3)
%  잔재를 발생시켜 enumerate 블록 전체 제거.]



%% -------------------------------------------------------------------
\subsection{\vrev{\revdel{이중채널 감쇠 효과와 P85.8 교차점}}}
\label{subsec:dual-channel-crossover}
%% -------------------------------------------------------------------

\fdel{본 연구의 가장 핵심적인 이론적 발견은 P85.8 분할 기준점(CDF 교차점과 일치)을
경계로 하는 이중채널 감쇠(dual-channel attenuation) 효과이다.
C3 조건에서 P85.8 기준으로 분할 회귀를 수행한 결과,
관측성 가중치 $k_O$의 회귀 계수가 Core 영역($\leq$ P85.8)의 $-0.798$에서
Tail 영역($>$ P85.8)의 $-0.415$로 $0.52$배 감쇠된다.}

\fdel{이 발견의 이론적 함의는 세 가지이다.
첫째, 극단적 불확실성 환경에서 관측성 투자의 \textbf{한계효용이 감소}한다.
Core 영역에서 $k_O$를 $0.1$ 증가시키면 $\eta$가 $0.080$ 감소하나,
Tail 영역에서는 동일한 변화가 $\eta$를 $0.042$만 감소시킨다.
이는 BCM 자원 배분에서 극단 시나리오일수록 관측성 투자의
비용-효과성이 감소하며 불확실성 관리가 우선되어야 함을 의미한다.
둘째, 시그모이드 비선형성과 파레토의 \textbf{상호작용 메커니즘}이
규명되었다. Tail 영역에서 큰 $U$ 값이 시그모이드의 포화 구간을 활성화하여
불확실성 채널이 지배적이 되며, $k_O$의 상대적 역할이 축소된다.
셋째, P85.8이 분할 회귀 기준점이자 CDF 교차점으로서 가우시안--파레토 분포 전환의
\textbf{구조적 경계}임이 다중 시드 검증을 통해 확인되어,
이 교차점이 특정 시드에 의존하지 않는 강건한 현상임이 입증되었다.}


  % --- 5.8.1 (이론적 기여) — 5.1.2와 3중 주제 중복으로 5-2 삭제 ---
  \begin{p5del}{현 5.8.1 이론적 기여 — 5.1.2의 1)/2)/3)과 동일 주제 중복}
  %% ===================================================================
\section{\vrev{학술적 기여}}
\label{sec:conclusion-contributions}
%% ===================================================================


%% -------------------------------------------------------------------
\subsection{\vrev{이론적 기여}}
\label{subsec:contrib-theoretical}
%% -------------------------------------------------------------------

\fdel{본 연구의 이론적 기여는 다음의 세 가지로 정리된다.}

\fdel{첫째, \textbf{업무연속성관리$\cdot$관측성$\cdot$시그모이드 바운딩의 통합적 이론 기반을
확립}하였다.
기존 BCM 연구에서 RTO는 BIA 시점에 고정되는 정적 매개변수로 취급되었으며,
관측성과 불확실성을 단일 수리 모델에서 통합한 선행 연구는 부재하였다.
본 연구는 제어 이론의 관측성 개념과 확률론적 불확실성 모형을
수정 시그모이드 함수라는 공통 기저 함수 위에서 결합함으로써,
BCM 분야에 새로운 이론적 기반을 제공한다.
특히 $\DRTO = R_0 + E_{O,\text{bnd}} + E_{U,\text{bnd}}$의 가산 구조는
각 효과의 독립적 분석과 결합 효과의 동시 포착을 가능하게 하여,
모델의 해석 가능성(interpretability)과 확장 가능성(extensibility)을
동시에 확보한다.}

\fdel{둘째, \textbf{꼬리 감쇠 현상을 발견하고 정량화}하였다.
P85.8 분할 기준을 경계로 관측성 가중치 $k_O$의 효과가
Core($-0.798$)에서 Tail($-0.415$)로 $0.52$배 감쇠되는 현상은,
극단 불확실성 환경에서 시그모이드 포화에 의해
관측성 채널의 조절 효과가 구조적으로 제한됨을 규명한 것이다.
이 발견은 BCM 자원 배분 전략에 새로운 분석적 기반을 제시하며,
극단 상황에서의 대비 전략이 관측성 투자만으로는 한계가 있으며
사전 자원 확보가 병행되어야 한다는 실무적 시사점을 제공한다.}

\fdel{셋째, \textbf{P85.8 CDF 교차점을 가우시안--파레토 분포 전환의 구조적 경계로 규명}하였다.
C2와 C3의 CDF가 P85.8($\eta \approx 2.42$)에서 교차하며,
이 교차점을 경계로 두 조건의 행태가 질적으로 전환된다.
P85.8 이하에서 두 분포는 사실상 구분 불가능하나,
P85.8 초과에서 C3의 $\eta$가 C2를 급격히 상회하여
CVaR$_{99}$ 기준 $+24.0\%$의 격차를 보인다.
이는 파레토 불확실성의 효과가 \textit{평균 수준의 비차별성}과
\textit{꼬리 영역의 극적 분기}로 구성되는 이중 구조를 가진다는 것을
실증적으로 규명한 결과이다.}


  \end{p5del}
\endgroup

% --- 5.2.2 방법론적 기여 ---
\subsection{\vrev{방법론적 기여}}
\label{subsec:academic-methodological}

\begingroup
  \renewcommand{\section}[1]{}
  \renewcommand{\subsection}[1]{}
  \renewcommand{\subsubsection}[1]{\paragraph{#1}}

  

%% ===================================================================
\section{\vrev{방법론적 시사점}}
\label{sec:methodological-implications}
%% ===================================================================

본 연구는 $\DRTO$ 프레임워크의 검증 과정에서 \jrev{네 가지} 방법론적 혁신을
달성하였다. \chg{이를 차례로 서술한다.}


%% -------------------------------------------------------------------
\subsection{\vrev{몬테카를로 시뮬레이션 방법론}}
\label{subsec:mc-methodology}
%% -------------------------------------------------------------------

\chg{첫째, 몬테카를로 시뮬레이션 방법론이다.} $n = 10{,}000$ 표본 $\times$ 4개 조건(C0--C3) $\times$ $10{,}000$개 마스터 시드의
총 $4 \times 10^8$회 $\DRTO$ 계산을 통해, 시뮬레이션 기반 BCM 모델 검증의
새로운 표준을 제시하였다.

\textbf{비중첩 대역 시드(non-overlapping band seeding)} 체계는 마스터 시드
$s = 42$로부터 변수별 시드를 $10{,}000$ 간격의 비중첩 대역\chg{(}$O_{\mathrm{raw}}$:
$[0\text{--}9999]$, $k_O$: $[20000\text{--}29999]$, $U^{(G)}$:
$[30000\text{--}39999]$, $U^{(P)}$: $[40000\text{--}49999]$\chg{)}으로
파생한다. \jrev{H6} 검증에서 6개 변수쌍 모두 $|r| < 0.02$(최대 $|r| = 0.016$)를
충족하여, 시뮬레이션의 완전한 재현성과 변수 간 통계적 독립성이 동시에 보장되었다.

이 체계의 방법론적 의의는 세 가지이다.
\chg{첫째,} 단일 시드로부터 다수의 독립적 난수열을 파생하는 표준화된 규약을 제공한다.
\chg{둘째,} $10{,}000$개 시드에 걸친 $10^8$회 다중 시드 검증(\jrev{H7})에서 Bonferroni 보정
$z$-검정 결과 45개 계수 중 기각이 $0$개로, 특정 시드에 대한 의존성을 완전히
배제한다. \chg{셋째,} 재현성 프로토콜이 명시되어 있어 후속 연구자의 독립적 검증이 가능하다.


%% -------------------------------------------------------------------
\subsection{\jrev{\texorpdfstring{\citet{lee2026drto}}{Lee and Cheung (2026)} $\kappa$ 매개변수의 본 연구 환경 적용 및 \frev{분포 환경 적합성} 시사 방법론}}
\label{subsec:kappa-optimization}
%% -------------------------------------------------------------------

\chg{둘째, $\kappa$ 매개변수의 분포 환경 적합성 시사 방법론이다.} \jrev{본 연구는 곡률 매개변수 $\kappa^{*} = 1.2$를 \chg{\citet{lee2026drto}}에서 다기준 종합점수($f_1$ 실무 유의성,
$f_2$ 비포화율, $f_3$ 효과 크기, $f_4$ 설명력)로 \chg{도출하여 출판된} 매개변수로
받아들이고, 본 연구의 새로운 분포 환경(\chg{C2 가우시안과 C3 파레토})에 그대로 적용하였다. 따라서 $\kappa^{*}$ 도출 자체는 본 연구의
독자적 기여가 아니며, \chg{\citet{lee2026drto}의} 격자 탐색 결과(C1 조건 기준
$F(\kappa^{*}) = 0.9997$)를 \textbf{전제 조건}으로 받아들인다.}

\jrev{본 연구의 방법론적 보완은 다음과 같다. \chg{\citet{lee2026drto}의}
4기준 종합점수 $F(\kappa) = f_1 \cdot f_2 \cdot f_3 \cdot f_4$는
C1 환경 기반으로 정의되어, 새로운 분포 환경(\chg{C2 가우시안과 C3 파레토})에서는
직접 격자 탐색에 부적합하다는 학술적 한계를 지닌다. 본 연구는 격자
탐색의 직접적 대안으로서, 9개 가설(\jrev{H2--H10}) PASS의 \textbf{통합 해석을
통한 \jrev{\frev{분포 환경 적합성}} 종합 명제 도출 방법론}을 제시하였다
(\secref{sec:ch4_hypothesis_summary} \jrev{\frev{분포 환경 적합성}} 종합 명제 참조).
구체적으로, \chg{i) 물리적 안전성(H2), ii) 이론적 일치(H3), iii) 통계적 설명력(H4·H5),
iv) 결과 안정성(H6·H7), v) 메커니즘 발현(H8), vi) 외부 검증(H9·H10)}의 6가지 차원에서 $\kappa = 1.2$의 \jrev{\frev{분포 환경 적합성}}이 집합적으로
\frev{시사된다}.}

\jrev{이러한 \textbf{다중 가설 PASS의 통합 해석} 방법론은 \chg{\citet{lee2026drto}}
매개변수의 새로운 환경 적용성을 \frev{시사}하는 보완적
방법론적 틀로서, 본 연구가 독자적으로 제시하는 기여이다. 본 접근법은
선행연구를 기반으로 분포 환경을 확장하는 후속 연구 일반에 적용 가능하며,
직접적 격자 탐색이 부적합한 환경에서의 매개변수 적합성 검증에 활용될 수 있다.}


%% -------------------------------------------------------------------
\subsection{\vrev{분할 회귀(Split Regression) 방법론}}
\label{subsec:split-regression}
%% -------------------------------------------------------------------

\chg{셋째,} P85.8 기준 분할 회귀는 파레토 분포의 체제 전환(regime change)을
포착하기 위한 방법론적 혁신이다. C3 전체 표본에서 \pdferr{2차 회귀}의
$R^2 = 0.858$에 그치는 잔여 비설명 분산($14.2\%$)을,
P85.8 기준으로 Core($\leq$ P85.8, $n = 8{,}580$)와 Tail($>$ P85.8, $n = 1{,}420$)로
분할함으로써 Core $R^2 = 0.999$, Tail $R^2 = 0.710$으로 개선하였다.

분할 회귀의 방법론적 기여는 다음과 같다.
\chg{i)} 단일 회귀 모형으로 포착할 수 없는 \textbf{구간별 이질적 구조}를
데이터 기반으로 식별하는 체계적 절차를 제시한다.
\chg{ii)} 분할 기준점(P85.8)의 결정이 사전적 이론(CDF 교차)에 근거하므로,
임의적 구간 분할의 자의성 문제를 해소한다.
\chg{iii)} 이 방법론은 BCM 분야에 국한되지 않고, 파레토 분포가 관여하는
금융 리스크 관리, 보험 계리, 자연 재해 모형 등에
범용적으로 확장 가능하다.

  
%% -------------------------------------------------------------------
\subsection{\vrev{방법론적 기여}}
\label{subsec:contrib-methodological}
%% -------------------------------------------------------------------

\fdel{방법론적 관점에서 본 연구는 다음의 세 가지 기여를 달성하였다.}

\fdel{첫째, \textbf{P85.8 기준 분할 회귀(split regression)를 파레토 분포 분석에 적용}하였다.
C3 전체 표본에서 2차교호작용 회귀의 $R^2 = 0.858$에 그치는 잔여 비설명 분산을,
P85.8 기준으로 Core($\leq$ P85.8)와 Tail($>$ P85.8)로 분할함으로써
Core $R^2 = 0.999$(0.9990), Tail $R^2 = 0.710$(0.7100)으로 Core 영역의
결정론적 구조를 완전히 포착하였다.
이 분할 전략은 파레토 분포의 체제 전환(regime change)을 포착하는
범용적 분석 방법론으로, BCM을 넘어 보험, 금융, 재난관리 등
파레토 현상이 핵심인 분야에 확장 적용될 수 있다.}

\fdel{둘째, \textbf{비중첩 대역 시드(non-overlapping band seeding) 체계}를 설계하였다.
마스터 시드 $s = 42$로부터 변수별 시드를 $10{,}000$ 간격의 비중첩 대역($O_{\mathrm{raw}}$: $[0\text{--}9999]$, $k_O$: $[20000\text{--}29999]$,
$U^{(G)}$: $[30000\text{--}39999]$,
$U^{(P)}$: $[40000\text{--}49999]$)으로 파생하여,
시뮬레이션의 완전한 재현성을 보장하면서도 변수 간 통계적 독립성을
유지한다(모든 변수 쌍 $|r| < 0.02$).}

\jrev{넷째}, \textbf{6단계 다층 검증(six-stage multi-layer validation) 프레임워크}를 구축하였다.
기반 검증(L0, Golden Compare), 시뮬레이션 검증, 회귀 검증, 강건성 검증, 이중채널 검증, 전문가 평가의
6개 단계(L0--L5)가 계층적 의존 구조를 형성하여 모델의 타당성을 다각적으로 보장한다.
정량적 시뮬레이션과 정성적 전문가 평가를 삼각검증(triangulation)함으로써
단일 방법론의 편향을 최소화하였으며, 이 프레임워크는
수리 모델 기반 연구의 검증 설계 표준으로 활용될 수 있다.



\endgroup

% --- 5.2.3 실증적 기여 (5-2단계 신규 작성) ---
\subsection{\vrev{실증적 기여}}
\label{subsec:academic-empirical}

\begin{p5add}
본 연구에서 실증적 기여는 시뮬레이션과 전문가 평가의 결합 검증을 통해 달성한
결과를 토대로 정리한다.

첫째, \textbf{\jrev{10개 연구 가설}(H1--\jrev{H10})이 사전 설정 기준을 모두 충족하여 모두
채택}되었다. 시뮬레이션 검증(H1--\jrev{H3}), 회귀 검증(\jrev{H4}--\jrev{H5}), 강건성 검증(\jrev{H6}--\jrev{H7}),
이중 채널 검증(\jrev{H8}), 전문가 평가(\jrev{H9}--\jrev{H10})의 6단계 다층 검증 프레임워크에서
계층적 의존 구조가 모두 만족됨으로써, $\DRTO$ 모델의 내부 정합성과
외부 타당성이 동시에 확인되었다.

둘째, \textbf{$10{,}000$개 마스터 시드에 걸친 $10^8$회 다중 시드 검증}을
통해 결과의 시드 의존성을 완전히 배제하였다. Bonferroni 보정 $z$-검정에서
45개 회귀 계수 중 기각이 0개로, 시뮬레이션 결과가 특정 난수열에 의존하지 않는
강건성을 정량적으로 입증하였다.

셋째, \textbf{BCM/DR 분야 전문가 5명에 의한 \chg{정량적 및 정성적} 혼합 검증}을 통해
실무 타당성을 확보하였다. 27문항 5점 리커트 척도에서 전체 평균
$4.46/5.0$($t(4) = 3.68$, $p = 0.011$, Cohen's $d = 1.65$)을 기록하였으며,
S-CVI/Ave $= 0.985$, Cronbach's $\alpha = 0.89$로 내용 타당도와 내적
일관성이 확인되었다. 심층 인터뷰에서는 5대 주제 15개 하위 주제 모두에서
$\geq 80\%$ 동의가 달성되어, 시뮬레이션 결과가 실무자의 경험적 관찰과
부합함을 삼각 검증하였다.

\jrev{넷째, \textbf{\frev{분포 환경 적합성} 종합 명제}를 도출하였다. \chg{\citet{lee2026drto}에서 격자
탐색으로 도출하여 출판한} $\kappa^{*} = 1.2$를 본 연구의
새로운 분포 환경에 적용한 결과,
9개 가설(\jrev{H2--H10})의 PASS가 \chg{i)} 물리적 안전성, \chg{ii)} 이론적 일치, \chg{iii)} 통계적
설명력, \chg{iv)} 결과 안정성, \chg{v)} 메커니즘 발현, \chg{vi)} 외부 검증의 6가지 측면에서
$\kappa = 1.2$의 \jrev{\frev{분포 환경 적합성}}을 집합적으로 \frev{시사한다} (\secref{sec:ch4_hypothesis_summary} 참조).
신규 격자 탐색을 직접 수행하지 않고도 다중 가설 PASS의 통합 해석을 통해
\chg{\citet{lee2026drto}의 매개변수의} 본 연구 환경 적용성을 학술적으로 정직하게 \frev{시사한 것이}
본 연구의 \chg{방법론적 및 실증적} 차별점이다.}
\end{p5add}
